\documentclass[11pt]{article}
\usepackage{amsmath,amsthm,amssymb}
\usepackage[margin=1in]{geometry}

\newtheorem{theorem}{Theorem}
\newtheorem{proposition}{Proposition}
\newtheorem{lemma}{Lemma}
\newtheorem{conjecture}{Conjecture}
\theoremstyle{definition}
\newtheorem{remark}{Remark}

\newcommand{\Z}{\mathbb{Z}}
\newcommand{\C}{\mathbb{C}}
\newcommand{\Zi}{\Z[i]}
\newcommand{\vpi}{v_\pi}

\title{Gap B: the two--row case of $G_\lambda(i)=0\iff\lambda=(2,2)$}
\author{Computational note (Clio compute agent)}
\date{2026-06-04}

\begin{document}
\maketitle

\section{Setup and the reduction chain}

For $\lambda\vdash n=2m$ write $G_\lambda(i)=\langle s_\lambda,\psi^m\rangle$ with
$\psi=h_2+i\,e_2$. The open conjecture is $G_\lambda(i)=0\iff\lambda=(2,2)$.
We treat two--row partitions $(a,b)$, $a\ge b\ge0$, $a+b=2m$.

Let
\[
  f(u)=u^2+(1+i)u+1\in\Zi[u],\qquad r_l:=[u^l]f(u)^m .
\]

\begin{proposition}[Reductions, verified]\label{prop:red}
The following hold and were checked exactly in $\Zi$ for all $n=2m\le 24$
(and the palindrome/formula identities for $m\le 12$):
\begin{enumerate}
\item In two variables $\psi(x_1,x_2)=x_1^2+x_2^2+(1+i)x_1x_2$ and
  $\psi^m=\sum_{l}r_l\,x_1^{l}x_2^{2m-l}$; consequently
  $G_{(a,b)}=r_a-r_{a+1}$ (Jacobi--Trudi $s_{(a,b)}=h_ah_b-h_{a+1}h_{b-1}$).
\item $f$ is palindromic, $u^2f(1/u)=f(u)$, hence $r_l=r_{2m-l}$.
\item Writing $b$ for the smaller part, $G_{(a,b)}=r_b-r_{b-1}$, $0\le b\le m$.
\item The closed form
  $r_l=\sum_{p=\max(0,l-m)}^{\lfloor l/2\rfloor}
        \dfrac{m!}{p!\,(m-l+p)!\,(l-2p)!}\,(1+i)^{\,l-2p}$
  agrees with $[u^l]f^m$.
\end{enumerate}
\end{proposition}

\noindent\emph{Verification.} An independent bivariate expansion of $\psi^m$
(\texttt{verify\_AB.py}) reproduces $G_{(a,b)}$ obtained from $r_a-r_{a+1}$; the
palindrome and the explicit formula were checked coefficient by coefficient. No
mismatch occurred. $\square$

Thus the two--row conjecture is equivalent to:
\begin{conjecture}\label{conj:main}
For the coefficient sequence $(r_0,\dots,r_m)$ of $f^m$ one has
$r_l\ne r_{l-1}$ for all $1\le l\le m$, \emph{except} $(m,l)=(2,2)$
(the shape $(2,2)$).
\end{conjecture}

\paragraph{Scale.} Conjecture~\ref{conj:main} was confirmed by exact $\Zi$
computation for all $m\le 120$ (i.e.\ $n\le 240$); the \emph{only} consecutive
equality is $(m,l)=(2,2)$, where $r_1=r_2=2+2i$.

\section{Three handles}

Set $d_l:=r_l-r_{l-1}$ and let $\pi=1+i$, the prime above $2$ in $\Zi$,
$\pi^2=2i$, $N(\pi)=2$; recall $\vpi(z)=v_2(N(z))$ with $N(a+bi)=a^2+b^2$.

\subsection{Modular reduction $\bmod\,\pi$ (a clean partial theorem)}

\begin{lemma}\label{lem:modpi}
$r_l\equiv\binom{2m}{l}\pmod\pi$ for all $l$. Equivalently, under
$\Zi/(\pi)\cong\mathbb{F}_2$, $\;a+bi\mapsto a+b$, one has
$r_l\equiv\binom{2m}{l}\bmod2$.
\end{lemma}
Verified for all $m\le 19$, every $l$.

\begin{proposition}[mod-$\pi$ proof, unconditional]\label{prop:notbotheven}
If \emph{not both} of $\binom{2m}{l},\binom{2m}{l-1}$ are even, then
$\vpi(d_l)=0$, hence $d_l\ne0$ and $r_l\ne r_{l-1}$.
\end{proposition}
\begin{proof}
By Lemma~\ref{lem:modpi}, $d_l\equiv\binom{2m}{l}-\binom{2m}{l-1}\pmod\pi$.
If exactly one of the two binomials is odd, the difference is odd, so
$d_l\not\equiv0\bmod\pi$ and $\vpi(d_l)=0$.
\emph{Both odd is impossible}: by Lucas, $\binom{2m}{l}$ is odd iff $l$ is a
binary submask of $2m$; since $2m$ is even its bit~$0$ is $0$, forcing $l$ even;
then $l-1$ is odd, has bit~$0$ equal to~$1$, and cannot be a submask of $2m$, so
$\binom{2m}{l-1}$ is even. Hence ``not both even'' $=$ ``exactly one odd''.
\end{proof}
Numerically: in every ``easy'' case $d_l$ is odd mod $\pi$ (no violation,
$m\le 99$). The leftover ``hard'' set is exactly
$\{\,l:\binom{2m}{l}\equiv\binom{2m}{l-1}\equiv0\ (2)\,\}$. For
$m=2^k-1$ ($m=3,7,15,31,63,\dots$) \emph{all} binomials $\binom{2m}{l}$ are odd,
so Proposition~\ref{prop:notbotheven} settles those $m$ entirely.

\subsection{$\pi$-adic depth is unbounded (why mod-$\pi$ alone cannot finish)}

On the hard (both-even) cases $\vpi(d_l)$ is positive and \emph{grows}: at
$m=2^k$ the maximal hard valuation equals $2k+2$
($m=4,8,16,32,64\Rightarrow 6,8,10,12,14$). Hence there is \emph{no} fixed level
$\pi^N$ at which $d_l$ is provably a unit for all $m$; a uniform $\pi$-adic
congruence proof is excluded. (Confirmed $m\le 79$.) This matches the previously
recorded obstruction ``$v_\pi$ depth unbounded $\Rightarrow$ no size bound''.

\subsection{Modulus monotonicity: the most promising handle}

Let $N_l:=|r_l|^2=\Re(r_l)^2+\Im(r_l)^2\in\Z_{\ge0}$ (cross-checked via
$N_l=r_l\cdot\overline{r_l}$ with $\overline{r_l}=[u^l]\bar f^m$,
$\bar f=u^2+(1-i)u+1$).

\begin{proposition}[verified to $m=200$]\label{prop:mono}
For every $m\ge1$ the real sequence $N_0,N_1,\dots,N_m$ is strictly increasing,
\[
  N_0<N_1<\cdots<N_m,
\]
with the \emph{single} exception $N_1=N_2\,(=8)$ at $m=2$.
\end{proposition}

This was checked exactly for all $m\le 200$ ($n\le400$): there are \emph{no}
decreases and exactly one tie, $(m,l)=(2,2)$.

\begin{theorem}[conditional]\label{thm:cond}
Proposition~\ref{prop:mono} implies Conjecture~\ref{conj:main}, hence Gap~B.
\end{theorem}
\begin{proof}
If $N_l\ne N_{l-1}$ then $|r_l|\ne|r_{l-1}|$ and in particular $r_l\ne r_{l-1}$.
The lone tie is $(m,l)=(2,2)$, the shape $(2,2)$, where indeed
$r_1=r_2=2+2i$.
\end{proof}

\paragraph{Why this is the right reduction.}
$N_l$ is a sequence of \emph{nonnegative integers}, free of the cancellation that
makes the complex $r_l$ hard. The binding constraint is the endpoint $l=m$
(the central coefficient $r_m=[u^m]f^m=\sum_p\binom{m}{p}\binom{m-p}{p}(1+i)^{m-2p}$,
the trinomial/Gegenbauer central value at argument $1+i$). For $m\ge3$ the minimum
of $N_l/N_{l-1}$ over $1\le l\le m$ is attained at $l=m$ and satisfies, with high
numerical precision,
\[
  \frac{N_m}{N_{m-1}}=1+\frac{3}{2m}+o\!\left(\tfrac1m\right),
\qquad
  m\Bigl(\tfrac{N_m}{N_{m-1}}-1\Bigr)\to \tfrac32 ,
\]
($m\,(\text{ratio}-1)=1.4994,1.4999$ at $m=200,300$). So the ratio is bounded
\emph{above} $1$ for every $m$, decreasing to $1$, with the worst case at the
center.

\section{Status and the remaining gap}

\begin{itemize}
\item \textbf{Reductions (A):} verified exactly, no mismatch ($n\le24$;
palindrome/formula $m\le12$).
\item \textbf{Scale (B):} Conjecture~\ref{conj:main} holds for all $m\le120$
($n\le240$); unique consecutive equality $(m,l)=(2,2)$.
\item \textbf{Proof (C):}
  \begin{itemize}
  \item mod-$\pi$ (Prop.~\ref{prop:notbotheven}) is an \emph{unconditional}
  proof for all $(m,l)$ outside the ``both-even'' set, and for all
  $m=2^k-1$ entirely.
  \item the $\pi$-adic valuation of $d_l$ is \emph{unbounded}, so no fixed-level
  congruence finishes the both-even cases (handle (ii) blocked).
  \item the modulus handle (iv) is the clean lever: Gap~B follows from the purely
  real statement that $|r_l|^2$ is strictly increasing on $0\le l\le m$ with the
  single exception at $m=2$ (Theorem~\ref{thm:cond}). The naive form of (iv)
  ($|r_l|$ monotone with no exception) is \emph{false} only at $(2,2)$, which is
  precisely the conjecture's exceptional shape.
  \end{itemize}
\end{itemize}

\noindent\textbf{Remaining gap.} Prove Proposition~\ref{prop:mono}:
$N_l<N_{l+1}$ for $0\le l<m$ (except $m=2$). A natural attack: the diagonal
generating function
\[
  \sum_{l\ge0}N_l\,z^l=\bigl[w^{2m}\bigr]\,Q(w)^m,\qquad
  Q(w)=f(w)\,\bigl(w^2+(1-i)zw+z^2\bigr),
\]
the central coefficient of a degree-$4$ polynomial power, plus a positivity/
log-concavity argument on the integer sequence $N_l$; the binding case $l=m$
has the explicit central-trinomial form above with the clean asymptotic ratio
$1+\tfrac{3}{2m}$.
\end{document}
